\documentclass{article}

% ICLR 2027 style — replace with official iclr2027_conference.sty when available
\usepackage{iclr2025_conference}
\usepackage{times}
\usepackage{hyperref}
\usepackage{url}
\usepackage{booktabs}
\usepackage{amsmath,amssymb}
\usepackage{graphicx}
\usepackage{xcolor}
\usepackage{multirow}
\usepackage{subcaption}

\title{Measuring the Value of Cognitive Sequences\\in LLM Research Agents}

\author{Anonymous}

\newcommand{\ie}{\textit{i.e.}}
\newcommand{\eg}{\textit{e.g.}}
\newcommand{\cf}{\textit{cf.}}

\iclrfinalcopy

\begin{document}

\maketitle

\begin{abstract}
Large language model (LLM) research agents combine heterogeneous cognitive operations---retrieval, hypothesis generation, reasoning, and falsification---to solve epistemic tasks.
Most architectures select operations greedily, assuming local value suffices to guide useful global reasoning.
We introduce a controlled epistemic benchmark enabling \emph{same-state counterfactual evaluation}: from a single epistemic state, we fork into alternative cognitive operations, apply matched continuation budgets, and compare downstream epistemic quality.
In controlled experiments across eight epistemic regimes ($N{=}56$ worlds), specific operation pairs exhibit super-additive temporal complementarity (hypothesis generation $\to$ retrieval: $+0.162$, 95\% CI $[0.135, 0.188]$), while repeating operations shows interference ($-0.050$, CI crosses zero).
Attack/falsification has zero standalone value but gains value after hypothesis formation ($+0.137$ for hypothesis $\to$ attack).
These structural effects persist when cognitive operations are accompanied by real LLM inference with two model substrates (Gemma~3 27B and Llama~3.2 11B).
We also report systematic negative findings: sequence superiority over primitives is not confirmed ($+0.003$), a learned adaptive controller fails (quality $0.323$ vs.\ best fixed $0.518$), and curated evidence tasks show negligible condition sensitivity (range $= 0.030$, $N{=}14$).
The benchmark, negative-results analysis, and evidence that the relevant unit of metacognitive analysis is operation-pair complementarity rather than general sequence superiority constitute our contributions.
\end{abstract}

%------------------------------------------------------------------
\section{Introduction}
\label{sec:intro}
%------------------------------------------------------------------

LLM-based research agents orchestrate qualitatively distinct \emph{cognitive operations}---evidence retrieval, hypothesis generation, causal reasoning, and hypothesis falsification---to solve complex epistemic tasks.
The central design question is: \emph{what cognitive operation should be performed next?}

Most architectures answer either with a fixed pipeline or with greedy scheduling, where the locally most valuable operation is selected at each step.
Both approaches assume that operation value can be assessed independently of temporal context.
This assumption may be wrong.
If \emph{temporal complementarity} exists---if the downstream value of operation $a_i$ depends on whether $a_j$ precedes or follows it---then greedy scheduling will systematically misallocate cognitive resources.

We investigate whether heterogeneous cognitive operations exhibit measurable temporal complementarity.
Our approach rests on a methodological contribution: a controlled epistemic benchmark that maintains latent ground truth, enabling \emph{same-state counterfactual evaluation}.
From a single epistemic state, we fork execution into alternative operations, apply matched continuation budgets, and compare downstream quality.
This isolates temporal-composition value from confounding state variation.

\paragraph{Summary of findings.}
Specific operation pairs show strong complementarity: hypothesis generation catalyzes retrieval ($+0.162$) and reasoning ($+0.137$), while attack has zero standalone value but gains value after hypothesis formation.
These structural effects persist when operations are accompanied by real LLM inference.
However, general sequence superiority over primitives is not confirmed, adaptive motif control fails, and effects do not transfer to curated evidence tasks.

\paragraph{Contributions.}
\begin{enumerate}
\item A controlled benchmark for counterfactual evaluation of cognitive operation sequences from matched epistemic states (\S\ref{sec:benchmark}).
\item Empirical evidence of temporal complementarity and interference among cognitive operations, with 95\% confidence intervals (\S\ref{sec:results_compl}).
\item Confirmation that complementarity effects persist under real LLM inference with two model substrates (\S\ref{sec:results_llm}).
\item Systematic negative results: adaptive control failure, non-transfer to curated evidence, and bounds on sequence superiority (\S\ref{sec:results_negative}).
\end{enumerate}

%------------------------------------------------------------------
\section{Related Work}
\label{sec:related}
%------------------------------------------------------------------

\paragraph{Adaptive test-time compute.}
Adaptive Computation Time \citep{graves2016act} and PonderNet \citep{banino2021pondernet} vary the \emph{amount} of computation.
We instead vary the \emph{type}: we ask whether qualitatively different operations compose non-additively.

\paragraph{Adaptive retrieval.}
Self-RAG \citep{asai2024selfrag} and FLARE \citep{jiang2023flare} decide \emph{when} to retrieve.
We measure the compositional value of retrieval relative to other operations.

\paragraph{Process reward models.}
Step-level rewards for mathematical reasoning \citep{lightman2024verify, wang2024mathshepherd} evaluate steps of the \emph{same kind}.
Our operations differ in kind---retrieval, hypothesis generation, attack, reasoning.

\paragraph{Reasoning trajectories.}
Tree of Thoughts \citep{yao2024tree}, Graph of Thoughts \citep{besta2024graph}, and LATS \citep{zhou2024lats} search over reasoning paths but do not isolate compositional value of heterogeneous operation types from matched states.

\paragraph{Metacognitive agents.}
Reflexion \citep{shinn2023reflexion} and Self-Refine \citep{madaan2023selfrefine} add self-monitoring.
Our negative results on adaptive control complement this literature.

\paragraph{Tool scheduling and temporal abstraction.}
Toolformer \citep{schick2023toolformer} and ART \citep{paranjape2023art} schedule tool use; the options framework \citep{sutton1999options} studies temporally extended actions.
Our work is analogous to options but in an epistemic rather than reward-maximizing context.

\paragraph{Distinction.}
To our knowledge, no prior work combines same-state counterfactual evaluation with explicit measurement of downstream complementarity among heterogeneous epistemic operations (Table~\ref{tab:prior_art}).

%------------------------------------------------------------------
\section{Problem Formulation}
\label{sec:formulation}
%------------------------------------------------------------------

\paragraph{Epistemic state.}
An epistemic state $E_t = (\mathcal{R}_t, \mathcal{H}_t, \mathcal{I}_t, c_t)$ captures retrieved evidence~$\mathcal{R}$, hypotheses~$\mathcal{H}$, inferred relations~$\mathcal{I}$, and cumulative cost~$c$.

\paragraph{Cognitive operations.}
We define four primitives: \textsc{Retrieve}~($a_\text{ret}$), \textsc{Hypothesize}~($a_\text{hyp}$), \textsc{Attack}~($a_\text{atk}$), and \textsc{Reason}~($a_\text{rsn}$).
Each transforms the epistemic state: $E_{t+1} = \tau(E_t, a_t)$.

\paragraph{Temporal complementarity.}
Two operations exhibit \emph{temporal complementarity} at state $E$ if their sequential value exceeds the additive baseline:
\begin{equation}
\label{eq:complementarity}
\text{Comp}(a_i, a_j \mid E) = Q(\tau(\tau(E, a_i), a_j)) - \tfrac{1}{2}\bigl[Q(\tau(E, a_i)) + Q(\tau(E, a_j))\bigr]
\end{equation}
where $Q$ is an evaluator-defined quality function.
When $\text{Comp} > 0$ (synergy), the sequence exceeds additive prediction; when $\text{Comp} < 0$ (interference), it underperforms.

\paragraph{Resource-fairness caveat.}
A 2-operation sequence consumes approximately twice the resources of a single operation.
Our complementarity measure therefore includes a resource confound: positive complementarity could partially reflect additional compute rather than temporal structure alone.
We report this limitation throughout.

\paragraph{Counterfactual fork.}
Given $E_t$, a counterfactual fork evaluates $G(S_A \mid E_t) - G(S_B \mid E_t)$ where both branches start from the same state.
This controls for state variation.

%------------------------------------------------------------------
\section{Benchmark: Epistemic World Simulator}
\label{sec:benchmark}
%------------------------------------------------------------------

The benchmark generates epistemic reasoning tasks with controlled latent structure.
Each \emph{world} contains: latent hypotheses (one true), an evidence corpus with differential support, source dependencies, and hidden variables affecting reliability.
The agent observes evidence but never the true hypothesis, causal graph, or quality function.

\paragraph{Epistemic regimes.}
Eight regimes (A--H) vary evidence informativeness, hypothesis discriminability, source independence, and deception structure.
Each regime instantiates multiple worlds via deterministic seeds.

\paragraph{Quality measurement.}
$Q(E)$ combines hypothesis correctness, evidence coverage, and calibration.
The function is evaluator-only: the agent never observes its score.

\paragraph{Same-state evaluation.}
The simulator supports event sourcing: any state can be snapshotted and restored, enabling counterfactual forks from identical starting conditions.
This is the methodological core of the paper.

%------------------------------------------------------------------
\section{Experiments and Results}
\label{sec:results}
%------------------------------------------------------------------

We organize experiments around five questions, summarized in Table~\ref{tab:claim_ladder}.

%------------------------------------------------------------------
\subsection{Temporal Complementarity and Interference}
\label{sec:results_compl}
%------------------------------------------------------------------

\paragraph{Setup.}
From matched epistemic states across 8 regimes ($N{=}56$ worlds, 7 seeds per regime), we execute all pairwise 2-operation sequences and compute $\text{Comp}(a_i, a_j)$ per Eq.~\ref{eq:complementarity}.

\paragraph{Results (Table~\ref{tab:complementarity}).}
Hypothesis generation is the catalytic operation: \textsc{Hyp}$\to$\textsc{Ret} yields complementarity $+0.162$ (95\% CI $[0.135, 0.188]$); \textsc{Hyp}$\to$\textsc{Atk} and \textsc{Hyp}$\to$\textsc{Rsn} each yield $+0.137$.
Repeating \textsc{Hyp}$\to$\textsc{Hyp} shows directional interference ($-0.050$, CI $[-0.136, +0.035]$; not significant at $\alpha{=}0.05$).

\textsc{Attack} and \textsc{Reason} have zero standalone quality gain.
They contribute value only after \textsc{Hypothesize}, creating a prerequisite-like structure: hypothesis generation ``unlocks'' downstream operations.

\paragraph{Interpretation.}
This pattern is partially expected: only \textsc{Hypothesize} creates candidate hypotheses, so \textsc{Attack} and \textsc{Reason} have nothing to evaluate without prior hypothesis generation.
The \emph{magnitude} of complementarity and its variation across regimes are the non-trivial findings, not the existence of prerequisite structure per se.
The 2-vs-1 operation resource confound means complementarity values include some contribution from additional compute.

%------------------------------------------------------------------
\subsection{Adaptive Control Failure}
\label{sec:results_adaptive}
%------------------------------------------------------------------

\paragraph{Adaptivity gap.}
A state-conditioned sequence oracle achieves quality $0.622$; the best global fixed sequence (B1\_extended) achieves $0.527$; gap $= 0.096$ ($N{=}56$).
Adaptive selection has theoretical room to improve.

\paragraph{Policy failure.}
A rule-based motif controller trained on state features achieves quality $0.323$---far below the best fixed strategy ($0.518$) and the oracle ($0.588$).
Policy regret ($0.264$) exceeds fixed regret ($0.070$) by $3.8\times$.
The controller produces only 2 distinct trajectory types.
This represents a complete failure of our adaptive-control approach.

\paragraph{Implications.}
The existence of an adaptivity gap does not imply that a learned controller can exploit it.
More expressive policy classes might succeed, but this remains a hypothesis.

%------------------------------------------------------------------
\subsection{LLM Execution}
\label{sec:results_llm}
%------------------------------------------------------------------

\paragraph{Models.}
Two LLM substrates pass a 10-operation capability gate ($\ge 0.73$ mean score, $N{=}24$ worlds each): Gemma~3 27B (Q4\_0 quantization) and Llama~3.2 11B (Q8\_0).
Both served via Ollama on a Mac Studio; temperature $= 0$.

\paragraph{Design.}
We execute 8 cognitive sequences across $N{=}16$ worlds (8 regimes $\times$ 2 seeds) with each LLM.
Quality is measured by the deterministic simulator---the LLM generates cognitive outputs but does \emph{not} affect the simulator's state transitions or quality evaluation.

\paragraph{Key design limitation.}
Because the simulator controls state transitions independently of LLM output, this experiment tests whether the structural complementarity effects \emph{persist when real LLMs perform the cognitive labor}, not whether LLM-generated cognition itself exhibits temporal complementarity.
Quality scores are identical across both models because they are determined by the simulator.

\paragraph{Results (Table~\ref{tab:replication}).}
\begin{itemize}
\item Complementarity: effect $+0.111$ (CI $[0.067, 0.156]$), consistent with V4 direction.
\item Attack timing: standalone attack $= 0.000$; after hypothesis+evidence: $0.276$.
\item Sequence superiority (B1\_extended vs.\ best primitive): $+0.003$, not significant.
\item Order effect (forward vs.\ reversed B1): $-0.050$, not significant.
\end{itemize}

The two strongest controlled effects (complementarity, attack timing) survive; weaker effects do not.

%------------------------------------------------------------------
\subsection{Curated Evidence Tasks}
\label{sec:results_real}
%------------------------------------------------------------------

\paragraph{Setup.}
14 curated evidence summaries across 6 academic domains, with evaluator-defined ground truth (keyword-heuristic scoring against author-specified correct hypotheses).
These are \emph{not} frozen real documents but literature-themed summaries with controlled structure.

\paragraph{Results (Table~\ref{tab:real_evidence}).}
Greedy primitive ($0.718$) achieves the highest mean score.
Condition range is $0.030$ across 5 strategies.
Cognitive sequence choice has negligible impact.

\paragraph{Interpretation.}
Multiple explanations are possible: (1) genuine non-transfer of complementarity effects; (2) insufficient task difficulty or strategy pressure in our curated packs; (3) LLMs may perform implicit hypothesis generation within single forward passes, reducing marginal value of explicit operations; (4) $N{=}14$ provides low statistical power (minimum detectable effect $\approx 0.05$).
We cannot distinguish among these with current evidence.

%------------------------------------------------------------------
\subsection{Negative Findings Summary}
\label{sec:results_negative}
%------------------------------------------------------------------

\begin{enumerate}
\item \textbf{Sequence superiority not confirmed under LLM:} B1\_extended ($0.277$) $\approx$ single hypothesis generation ($0.273$).
\item \textbf{Adaptive motif control failed:} Policy ($0.323$) $\ll$ best fixed ($0.518$).
\item \textbf{Curated-evidence transfer not detected:} Condition range $= 0.030$.
\item \textbf{Interference not significant:} $-0.050$, CI crosses zero.
\item \textbf{Full REE $<$ simple baselines:} The most complex architecture ($0.356$) underperformed a fixed sequence ($0.594$).
\end{enumerate}

These findings bound the scope of cognitive sequencing: temporal complementarity is operation-pair-specific rather than a general ``more steps $=$ better'' principle.

%------------------------------------------------------------------
\section{Discussion}
\label{sec:discussion}
%------------------------------------------------------------------

\paragraph{The unit of metacognitive control.}
Our data suggest that primitive step-level value can miss temporal interactions, but the failed motif controller shows that identifying sequence structure does not automatically yield a successful adaptive scheduler.
This tension---real complementarity, unsuccessful exploitation---is a central open problem for agent design.

\paragraph{Prerequisites vs.\ complementarity.}
A skeptical reading of our complementarity results is that they merely reflect prerequisite structure: \textsc{Attack} needs hypotheses; \textsc{Reason} needs hypotheses.
We partially agree---the \emph{existence} of these interactions is not surprising.
What is measured is their \emph{magnitude}, \emph{state dependence}, and \emph{persistence under LLM execution}.

\paragraph{Why transfer may fail.}
Modern LLMs may perform implicit hypothesis generation and reasoning within single forward passes.
Explicit cognitive operations may partially duplicate internal computation, reducing their marginal value on real tasks.

%------------------------------------------------------------------
\section{Limitations}
\label{sec:limitations}
%------------------------------------------------------------------

\begin{enumerate}
\item \textbf{Resource confound.} 2-operation sequences use $\sim$2$\times$ the tokens of single operations.
\item \textbf{Benchmark circularity risk.} The simulator was designed by the experimenters; findings may not generalize beyond it.
\item \textbf{Only two LLMs tested.} Both served from the same backend with the same prompts.
\item \textbf{LLM does not affect quality measurement.} Phase 26 quality comes from the simulator, not LLM output.
\item \textbf{Curated evidence, not real documents.} Phase 27 uses author-created summaries, not frozen source texts.
\item \textbf{Small sample sizes.} $N{=}16$ (LLM), $N{=}14$ (curated evidence).
\item \textbf{No formal hypothesis testing.} We report CIs but no formal tests; regime clustering likely makes CIs anti-conservative.
\item \textbf{No live-web evaluation.}
\item \textbf{Metric dependence.} Different quality scalarizations might yield different complementarity patterns.
\item \textbf{Adaptive policy tested only one controller class.}
\end{enumerate}

%------------------------------------------------------------------
\section{Conclusion}
\label{sec:conclusion}
%------------------------------------------------------------------

We presented a controlled benchmark for measuring temporal-composition value of cognitive operations.
Specific operation pairs exhibit measurable complementarity ($+0.162$) and prerequisite structure, and these effects persist under real LLM execution.
Equally important, general sequence superiority, adaptive control, and curated-evidence transfer are not supported.
The correct unit of metacognitive analysis appears to be the operation pair, not the full sequence or the individual primitive.

%------------------------------------------------------------------
% TABLES
%------------------------------------------------------------------

\begin{table}[t]
\centering
\caption{Pairwise complementarity matrix (Eq.~\ref{eq:complementarity}; $N{=}56$ worlds, 8 regimes). Bold: significant at 95\% CI.}
\label{tab:complementarity}
\small
\begin{tabular}{lcccc}
\toprule
First $\to$ Second & \textsc{Ret} & \textsc{Hyp} & \textsc{Atk} & \textsc{Rsn} \\
\midrule
\textsc{Retrieve}    & $-0.002$ & $+0.078$ & $+0.025$ & $+0.025$ \\
\textsc{Hypothesize} & $\mathbf{+0.162}$ & $-0.050$ & $\mathbf{+0.137}$ & $\mathbf{+0.137}$ \\
\textsc{Attack}      & $+0.025$ & $\mathbf{+0.137}$ & $0.000$ & $0.000$ \\
\textsc{Reason}      & $+0.025$ & $\mathbf{+0.137}$ & $0.000$ & $0.000$ \\
\bottomrule
\end{tabular}
\vspace{-2mm}
\end{table}

\begin{table}[t]
\centering
\caption{LLM execution results ($N{=}16$ worlds per model, 128 total evaluations). Quality from deterministic simulator.}
\label{tab:replication}
\small
\begin{tabular}{lccl}
\toprule
Test & Effect & 95\% CI & Verdict \\
\midrule
Complementarity & $+0.111$ & $[0.067, 0.156]$ & Consistent \\
Attack timing & $+0.276$ & $[0.179, 0.373]$ & Consistent \\
Sequence $>$ primitive & $+0.003$ & $[-0.160, 0.166]$ & Not confirmed \\
Order effect & $-0.050$ & $[-0.128, 0.028]$ & Not confirmed \\
\bottomrule
\end{tabular}
\vspace{-2mm}
\end{table}

\begin{table}[t]
\centering
\caption{Curated evidence results ($N{=}14$ packs, Gemma~3 27B). Condition range $= 0.030$.}
\label{tab:real_evidence}
\small
\begin{tabular}{lcc}
\toprule
Condition & Mean quality & Std \\
\midrule
Greedy primitive & $\mathbf{0.718}$ & $0.074$ \\
Direct & $0.696$ & $0.058$ \\
B1\_extended & $0.691$ & $0.112$ \\
Full explore & $0.688$ & $0.112$ \\
\bottomrule
\end{tabular}
\vspace{-2mm}
\end{table}

\begin{table}[t]
\centering
\caption{Cross-level claim ladder.}
\label{tab:claim_ladder}
\small
\begin{tabular}{lcccl}
\toprule
Finding & Simulator & +LLM & +Curated & Level \\
\midrule
Complementarity & \checkmark & \checkmark & --- & 2 \\
Attack timing & \checkmark & \checkmark & n/a & 2 \\
Interference & $\sim$ & $\sim$ & --- & 1 \\
Seq.\ $>$ prim. & \checkmark & --- & --- & 1 \\
Order effects & \checkmark & --- & --- & 1 \\
Adaptive ctrl. & gap exists & fails & n/a & 0 \\
\bottomrule
\end{tabular}
\vspace{1mm}
{\footnotesize \checkmark = supported, --- = not confirmed, $\sim$ = directional (CI crosses 0). \\
Level 2: simulator + LLM. Level 1: simulator only. Level 0: unsupported.}
\vspace{-3mm}
\end{table}

\begin{table}[t]
\centering
\caption{Comparison with related work. \checkmark = addressed, $\sim$ = partial, blank = not addressed.}
\label{tab:prior_art}
\small
\setlength{\tabcolsep}{3pt}
\begin{tabular}{lccccccc}
\toprule
& \rotatebox{70}{Adaptive} & \rotatebox{70}{Retrieval} & \rotatebox{70}{Op value} & \rotatebox{70}{Seq.\ value} & \rotatebox{70}{Same-state} & \rotatebox{70}{Compl.} & \rotatebox{70}{Hetero.\ ops} \\
\midrule
Self-RAG & $\sim$ & \checkmark & & & & & \\
ToT & \checkmark & & $\sim$ & $\sim$ & & & \\
LATS & \checkmark & & & \checkmark & & & \\
Reflexion & \checkmark & & & $\sim$ & & & $\sim$ \\
PRM & & & \checkmark & & & & \\
Options & \checkmark & & & \checkmark & & & \\
\textbf{This work} & \checkmark & \checkmark & \checkmark & \checkmark & \checkmark & \checkmark & \checkmark \\
\bottomrule
\end{tabular}
\vspace{-3mm}
\end{table}

%------------------------------------------------------------------
\bibliography{references}
\bibliographystyle{iclr2025_conference}

%------------------------------------------------------------------
% APPENDIX
%------------------------------------------------------------------
\appendix

\section{Reproducibility Statement}
\label{app:repro}

All experiments are deterministic given seeds and Python~3.11+.
Simulator experiments require no external dependencies.
LLM experiments used Gemma~3 27B QAT (digest \texttt{29eb0b9a...}) and Llama~3.2 11B Q8\_0 (digest \texttt{a5b7471a...}) via Ollama on a Mac Studio; temperature~$=0$.
Model quantization, generation parameters, and prompt templates are documented in the supplementary material.
Raw result artifacts (JSONL/JSON) are available for all experiments.

\section{AI Use Statement}
\label{app:ai_use}

This research used large language models in two capacities: (1)~as experimental subjects---Gemma~3 27B and Llama~3.2 11B were evaluated as cognitive-operation executors in the LLM replication experiments; (2)~as research assistants---an LLM assistant was used during code development, experimental design iteration, and manuscript drafting.
No LLM is listed as an author.
All scientific claims, experimental designs, and analyses were reviewed and validated by human researchers.

\section{Additional Experimental Details}
\label{app:details}

\paragraph{Capability gate.}
Both models were evaluated on 10 operations $\times$ 24 worlds.
Key scores: hypothesis generation (Gemma: $0.750$, Llama: $0.733$), reasoning (Gemma: $0.929$, Llama: $0.950$), all others $\ge 0.800$.
Gate threshold: $0.30$ per core operation.

\paragraph{Adaptive control details.}
The policy used transparent rule-based features from epistemic state (evidence count, hypothesis count, regime indicators).
Selection was among 5 motifs (explore, discriminate, full\_explore, integrate, evidence\_heavy).
$N{=}56$ training worlds; no held-out validation split for policy selection.

\paragraph{Curated evidence pack construction.}
14 packs were authored programmatically with literature-themed content referencing real academic debates (dual-process theory, emergent abilities, minimum wage effects, etc.).
Ground truth was defined by the pack author, not independently annotated.
Quality scoring uses keyword heuristics against author-specified correct hypotheses.
Packs are \emph{not} frozen real documents.

\paragraph{Budget fairness.}
2-operation sequences consume approximately $2\times$ the tokens of single operations (mean: 485 vs.\ 232 tokens).
Longer experimental sequences (B1\_extended, 7~ops) consume approximately $10\times$ single-operation budgets.
Complementarity measurements include this resource confound; see \S\ref{sec:formulation}.

\end{document}
